% Irodov Part III §3.2, original 78--88, PDF pp.111--112,
% printed pp.102--103; Fig 3.10--3.14 checked against source images.
\electroconductproblem{78}{102}
玻璃与真空分界面上的点 A
（见图3.10），真空侧电场强度
$E_0=10.0\,\mathrm{V/m}$，
与从玻璃指向真空的法线夹角
$\alpha_0=30^\circ$。
求玻璃内近边界处的场强 $E$、
与法线夹角 $\alpha$、
以及玻璃的束缚表面电荷密度。
\exsolution
设玻璃相对介电常数为 $\varepsilon_r$。
界面无自由面电荷时，
法向 $\bm D$ 连续、切向 $\bm E$ 连续：
\[
E_n=\frac{E_0\cos\alpha_0}{\varepsilon_r},
\qquad E_t=E_0\sin\alpha_0.
\]
因此
\[
\boxed{E=E_0\sqrt{
\frac{\cos^2\alpha_0}{\varepsilon_r^2}
+\sin^2\alpha_0}},
\qquad
\boxed{\tan\alpha=\varepsilon_r\tan\alpha_0}.
\]
玻璃极化
$\bm P=\varepsilon_0(\varepsilon_r-1)\bm E$，
故面束缚电荷
\[
\boxed{\sigma_b=
\varepsilon_0\frac{\varepsilon_r-1}{\varepsilon_r}
E_0\cos\alpha_0}.
\]
取原书数值答案所用的玻璃介电常数
$\varepsilon_r\simeq6$，
得到
\[
\boxed{E\simeq5.2\,\mathrm{V/m}},\quad
\boxed{\alpha\simeq74^\circ},\quad
\boxed{\sigma_b\simeq64\,\mathrm{pC/m^2}}.
\]
\excheck 玻璃的 $\varepsilon_r$ 未在本题直接列数值，
因此数值须说明采用的常用玻璃介电常数，
公式不依赖这一选择。

\electroconductproblem{79}{102}
一均匀各向同性电介质半空间的表面为平面，
其相对介电常数为 $\varepsilon_r$。
真空侧近表面的均匀电场
$\bm E_0$ 与法线成 $\theta$ 角
（图3.11）。
假设两侧场分别均匀，求：
(a) 以界面上点为球心、半径 $R$ 的球面的
$\bm E$ 通量；
(b) 沿图中与 $\bm E_0$ 平行的法截面内、
跨越界面的长为 $\ell$ 的狭矩形闭路的
$\bm D$ 环流。
\exsolution
无自由面电荷时
$E_{{\rm diel},n}=E_0\cos\theta/\varepsilon_r$，
$E_{{\rm diel},t}=E_0\sin\theta$。
(a) 球面的一半处于真空、一半在电介质。
每半球均匀场的通量由投影面积
$\pi R^2$ 给出，
所以
\[
\Phi_E=
\pi R^2E_0\cos\theta-
\pi R^2\frac{E_0\cos\theta}{\varepsilon_r}.
\]
即
\[
\boxed{\Phi_E=
\frac{\varepsilon_r-1}{\varepsilon_r}
\pi R^2 E_0\cos\theta}.
\]
(b) 让狭矩形两个长边分别位于界面的
真空侧和介质侧，
并按图3.11给定方向取闭路。
切向电场连续，
但切向位移分别为
$D_{0,t}=\varepsilon_0E_0\sin\theta$、
$D_{{\rm diel},t}=\varepsilon_0\varepsilon_rE_0\sin\theta$。
跨界面短边在宽度趋零时不贡献，
所以
\[
\boxed{\oint_\Gamma\bm D\cdot\mathrm d\bm l
=-\varepsilon_0(\varepsilon_r-1)
\ell E_0\sin\theta}.
\]
\excheck $\bm E$ 满足无旋性，
但 $\bm D$ 穿越不同介电常数的边界时
其切向分量一般不连续，
故 $\oint\bm D\cdot\mathrm d\bm l$
可以非零。

\electroconductproblem{80}{102}
一块沿 $x$ 轴法向、厚度为 $2d$
的无限大均匀电介质平板，
相对介电常数为 $\varepsilon_r$，
在 $|x|<d$ 内存在均匀\emph{自由体电荷密度}
$\rho_f$。
设中心电势为零，求：
(a) 电场 $E_x(x)$、电势 $\varphi(x)$；
(b) 束缚电荷的体密度和表面密度。
\exsolution
系统关于 $x=0$ 对称，
$\bm D$ 在介质内部满足
$\mathrm dD_x/\mathrm dx=\rho_f$，
由 $D_x(0)=0$，
$D_x=\rho_fx$。
在外侧
$D_x=\rho_fd\,\operatorname{sgn}x$。
因此
\[
\boxed{E_x(x)=
\begin{cases}
\dfrac{\rho_f x}{\varepsilon_0\varepsilon_r},
&|x|\le d,\\[2mm]
\dfrac{\rho_fd}{\varepsilon_0}\operatorname{sgn}x,
&|x|\ge d.
\end{cases}}
\]
用 $E_x=-\varphi'(x)$，
并取 $\varphi(0)=0$、电势连续，
\[
\boxed{\varphi(x)=
\begin{cases}
-\dfrac{\rho_fx^2}{2\varepsilon_0\varepsilon_r},
&|x|\le d,\\[2mm]
-\dfrac{\rho_fd^2}{2\varepsilon_0\varepsilon_r}
-\dfrac{\rho_fd}{\varepsilon_0}(|x|-d),
&|x|\ge d.
\end{cases}}
\]
极化
$\bm P=(\varepsilon_r-1)\bm D/\varepsilon_r$。
内部束缚体密度
\[
\boxed{\rho'_b=-\nabla\cdot\bm P
=-\rho_f\frac{\varepsilon_r-1}{\varepsilon_r}}.
\]
两侧束缚面电荷均同号且大小相等：
\[
\boxed{\sigma'_b(\pm d)=
+\rho_fd\frac{\varepsilon_r-1}{\varepsilon_r}}.
\]
\excheck 电场为\emph{奇函数}，
电势为\emph{偶函数}；
两侧束缚面电荷与内部束缚体电荷
总量每单位面积相消：
$2\sigma'_b+2d\rho'_b=0$。

\electroconductproblem{81}{102}
半径 $R$ 的均匀各向同性电介质实心球
相对介电常数为 $\varepsilon_r$，
内部自由体电荷密度
$\rho_f>0$ 均匀分布。
求 (a) 内外电场与电势的半径函数；
(b) 束缚电荷体密度与表面密度；
说明图线形状。
\exsolution
高斯定律对位移矢量：
$D_r(r)=\rho_fr/3$（$r<R$），
$D_r(r)=\rho_fR^3/(3r^2)$（$r>R$）。
因此
\[
\boxed{E(r)=
\begin{cases}
\dfrac{\rho_fr}{3\varepsilon_0\varepsilon_r},
&r<R,\\[2mm]
\dfrac{\rho_fR^3}{3\varepsilon_0r^2},
&r>R.
\end{cases}}
\]
取无穷远电势为零，
球面
$\varphi(R)=\rho_fR^2/(3\varepsilon_0)$，
由内场积分得到
\[
\boxed{\varphi(r)=
\begin{cases}
\dfrac{\rho_fR^2}{3\varepsilon_0}
+\dfrac{\rho_f(R^2-r^2)}
{6\varepsilon_0\varepsilon_r},&r\le R,\\[2mm]
\dfrac{\rho_fR^3}{3\varepsilon_0r},&r\ge R.
\end{cases}}
\]
内部极化
$\bm P=(\varepsilon_r-1)\rho_f\bm r/(3\varepsilon_r)$，
所以
\[
\boxed{\rho'_b=
-\rho_f\frac{\varepsilon_r-1}{\varepsilon_r}},
\quad
\boxed{\sigma'_b=
\frac{\varepsilon_r-1}{\varepsilon_r}
\frac{\rho_fR}{3}}.
\]
\excheck $E(r)$ 在球内线性增加，
在 $r=R$ 因介电常数改变而跳跃，
球外按 $1/r^2$ 衰减；
电势连续，内段为抛物线、外段为 $1/r$。
全部束缚体、面电荷代数和为零。

\electroconductproblem{82}{102}
半径 $R$、厚度 $d\ll R$ 的薄圆盘形电介质，
具有均匀固定极化强度 $\bm P$，
其方向\emph{位于圆盘平面内}。
求圆盘中心电场。
\exsolution
取 $\bm P=P\bm e_x$。
上下两个平面因其法线与 $\bm P$ 垂直，
束缚面电荷为零；
只有侧面 $r=R$ 带束缚电荷
$\sigma_b=P\cos\phi$。
在圆盘中心，侧面元
$\mathrm dS=R\,\mathrm d\phi\,\mathrm dz$
产生
$\mathrm d\bm E=
-k\sigma_b(R\bm e_r+z\bm e_z)
\,\mathrm dS/(R^2+z^2)^{3/2}$。
积分侧壁
$z\in[-d/2,d/2]$ 及角度，
由
$\int_0^{2\pi}\cos^2\phi\,\mathrm d\phi=\pi$，
得
\[
\bm E(0)=
-\frac{PR^2}{4\varepsilon_0}
\left(\int_{-d/2}^{d/2}
\frac{\mathrm dz}{(R^2+z^2)^{3/2}}\right)\bm e_x.
\]
在 $d\ll R$ 时，
括号内约为 $d/R^3$，
故
\[
\boxed{\bm E(0)\simeq
-\frac{d}{4\varepsilon_0R}\bm P}.
\]
\excheck 圆盘极化平行盘面，
故只有侧壁带束缚电荷；
若误当成盘面法向极化，将得到完全不同的场。

\electroconductproblem{83}{102--103}
无自由电荷的无限大电介质平板，
厚度为 $2d$，取中面为 $x=0$。
其极化分布
$\bm P(x)=\bm P_0(1-x^2/d^2)$，
$\bm P_0$ 垂直平板。
求板内电场与两表面间电势差。
\exsolution
设 $\bm P_0=P_0\bm e_x$。
没有自由电荷，
$\mathrm dD_x/\mathrm dx=0$。
且极化在 $|x|=d$ 处为零，
无外加电场的孤立板在外侧场为零，
故 $D_x=0$，即
$\varepsilon_0E_x+P_x=0$。
因此
\[
\boxed{\bm E(x)=
-\frac{\bm P_0}{\varepsilon_0}
\left(1-\frac{x^2}{d^2}\right)},\qquad |x|\le d.
\]
两表面间电势差定义为
$\varphi(d)-\varphi(-d)$，
有
\[
\Delta\varphi=
-\int_{-d}^dE_x\,\mathrm dx
=\frac{P_0}{\varepsilon_0}
\int_{-d}^d(1-x^2/d^2)\,\mathrm dx,
\]
故
\[
\boxed{\Delta\varphi=\frac{4dP_0}{3\varepsilon_0}}.
\]
\excheck 电介质平板外侧场为零是此\emph{无外加场}
的边界条件；否则 $D$ 可能带另一常数。

\electroconductproblem{84}{103}
两板间初始为空气的平行板电容器，
间距 $d$，初电场 $E_0$。
现按图3.12把\emph{沿场强方向的半层厚度}
替换为相对介电常数 $\varepsilon_r$ 的材料，
形成两个各厚 $d/2$ 的\emph{串联层}。
求气隙部分1和介质部分2的
$\bm E,\bm D$ 模：
(a) 极板电压不变；(b) 极板自由电荷不变。
\exsolution
界面没有自由面电荷，故
$D_1=D_2=D$。
两层法向场
$E_1=D/\varepsilon_0$、
$E_2=D/(\varepsilon_0\varepsilon_r)$。
(a) 电压固定为初值 $V_0=E_0d$，
\[
V_0=\frac d2(E_1+E_2)
=\frac{dD}{2\varepsilon_0}
(1+1/\varepsilon_r).
\]
所以
\[
\boxed{E_1=\frac{2\varepsilon_r}{\varepsilon_r+1}E_0},
\qquad
\boxed{E_2=\frac2{\varepsilon_r+1}E_0},
\quad
\boxed{D_1=D_2=
\frac{2\varepsilon_r}{\varepsilon_r+1}\varepsilon_0E_0}.
\]
(b) 极板自由电荷不变，
$D=D_0=\varepsilon_0E_0$，因此
\[
\boxed{E_1=E_0,\quad E_2=E_0/\varepsilon_r},
\qquad
\boxed{D_1=D_2=\varepsilon_0E_0}.
\]
\excheck 图3.12 是\emph{沿电场方向的串联半层}，
与第85题的并联半面积构型不同。

\electroconductproblem{85}{103}
把前题的电介质改按图3.13布置，
占电容器极板面积的一半，
且每块区域均跨越\emph{全板间距}，
形成空气区1与介质区2的\emph{并联}。
其余条件相同。求 (a) 电压不变；
(b) 极板电荷不变时各区 $E,D$。
\exsolution
忽略边缘效应时，
两区域上下极板相同，
因此具有相同电压 $V$，
$E_1=E_2=V/d$。
(a) $V=V_0=E_0d$，
故
\[
\boxed{E_1=E_2=E_0},\qquad
\boxed{D_1=\varepsilon_0E_0,\quad
D_2=\varepsilon_r\varepsilon_0E_0}.
\]
(b) 极板总面积为 $S$，
原自由电荷
$Q_0=\varepsilon_0E_0S$。
新电容为两个面积 $S/2$ 的电容并联：
\[
C_{\rm new}=\frac{\varepsilon_0S}{2d}
(1+\varepsilon_r).
\]
因此
$V=Q_0/C_{\rm new}
=2E_0d/(1+\varepsilon_r)$，
\[
\boxed{E_1=E_2=\frac{2E_0}{1+\varepsilon_r}},
\]
\[
\boxed{D_1=\frac{2\varepsilon_0E_0}{1+\varepsilon_r},
\quad
D_2=\frac{2\varepsilon_r\varepsilon_0E_0}
{1+\varepsilon_r}}.
\]
\excheck 固定\emph{总电荷}不表示两半极板
各自的自由电荷密度保持原值；
并联两区的自由电荷会重新分配。

\electroconductproblem{86}{103}
一个同心球形电容器，
内外导体球壳同心、内导体带自由电荷 $q$。
两极板之间的空间按图3.14
沿穿过球心的平面划为两半：
一半是真空，一半填相对介电常数
$\varepsilon_r$ 的均匀材料。
求极板之间的场强模 $E(r)$。
\exsolution
球形等势导体的边界迫使电场沿径向。
两半空间共用同一径向电势梯度，
故两侧 $E_r(r)$ 相同，
而 $D$ 不同：
$D_{\rm vac}=\varepsilon_0 E_r$，
$D_{\rm diel}=\varepsilon_0\varepsilon_r E_r$。
对半径 $r$ 的球面取高斯通量，
真空和介质各占半球面积 $2\pi r^2$：
\[
q=\oint\bm D\cdot\mathrm d\bm S
=2\pi r^2\varepsilon_0E_r+
2\pi r^2\varepsilon_r\varepsilon_0E_r.
\]
因此
\[
\boxed{E(r)=
\frac{|q|}{2\pi\varepsilon_0(1+\varepsilon_r)r^2}}.
\]
\excheck 两半分区\emph{并联}，不能把介质
误当成径向串联层；
当 $\varepsilon_r=1$ 恢复
$|q|/(4\pi\varepsilon_0r^2)$。

\electroconductproblem{87}{103}
两个带同号等量电荷的相同小球，
由等长绝缘线悬挂于同一点。
将周围介质由空气换成煤油后，
两线间张角\emph{保持不变}。
求小球材料密度 $\rho$，
煤油密度记为 $\rho_0$，
相对介电常数为 $\varepsilon_r$。
\exsolution
在空气中静电排斥力为 $F_e$，
球的有效重力约为 $mg$。
在煤油中，同样球间距离下
库仑力降为 $F_e/\varepsilon_r$；
与此同时浮力使有效重力变为
\[
mg-\rho_0Vg=(\rho-\rho_0)Vg
=mg\left(1-\frac{\rho_0}{\rho}\right).
\]
张角不变要求
\[
\frac{F_e}{mg}
=\frac{F_e/\varepsilon_r}
{mg(1-\rho_0/\rho)}.
\]
约去公共因子，
$1-\rho_0/\rho=1/\varepsilon_r$，
得
\[
\boxed{\rho=
\frac{\varepsilon_r}{\varepsilon_r-1}\rho_0}.
\]
取煤油常用近似
$\varepsilon_r\simeq2.0$、
$\rho_0\simeq0.80\,\mathrm{g/cm^3}$，
\[
\boxed{\rho\simeq1.6\,\mathrm{g/cm^3}}.
\]
\excheck 不能只考虑介电常数使库仑力减小，
必须同时计入煤油对小球的浮力。

\electroconductproblem{88}{103}
一个半径
$R=3.0\,\mathrm{cm}$ 的均匀各向同性电介质球，
相对介电常数为
$\varepsilon_r=5.00$，
球内部产生均匀电场
$E=100\,\mathrm{V/m}$。
求：(a) 最大束缚面电荷密度；
(b) 同号束缚电荷的总量。
\exsolution
线性介质的极化矢量为
$\bm P=\varepsilon_0(\varepsilon_r-1)\bm E$。
均匀极化球的表面束缚电荷
$\sigma_b(\theta)=\bm P\cdot\bm n
=P\cos\theta$。
故最大面密度位于极化轴两极：
\[
\boxed{\sigma_{b,\max}=
\varepsilon_0(\varepsilon_r-1)E}.
\]
代入数值：
\[
\boxed{\sigma_{b,\max}\simeq
(8.854\times10^{-12})(4)(100)
=3.5\,\mathrm{nC/m^2}}.
\]
同号电荷全部位于一半球面，
其总量是
\[
Q'_+=\int_0^{2\pi}\int_0^{\pi/2}
P\cos\theta R^2\sin\theta
\,\mathrm d\theta\,\mathrm d\phi
=\pi R^2P.
\]
因此
\[
\boxed{Q'_+=\pi R^2
\varepsilon_0(\varepsilon_r-1)E
\simeq1.0\times10^{-11}\,\mathrm C
=10\,\mathrm{pC}}.
\]
\excheck 负束缚电荷总量为 $-Q'_+$，
所以中性电介质球的束缚电荷代数和为零。
